\documentclass[11pt,a4paper,scheme=plain,fontset=fandol,no-math]{ctexart}
\usepackage[margin=25mm]{geometry}
\usepackage{mathtools}
\usepackage{eqnlines}
\usepackage[libertinus]{termes-otf}
\usepackage[restoremathleading=true]{zhlineskip}
\AtBeginDocument{\fontsize{11pt}{14.5pt}\selectfont}
\newcommand{\norm}[1]{\left\lVert #1 \right\rVert}
\newcommand{\E}{\mathbb{E}}
\begin{document}
扩散模型的验证指标涉及期望、平方根、嵌套函数与范数。

\begin{equation}\label{eq:diffusion}
\E\left[ \norm{ x_{t+1} - x_t }^2 \right] = \sqrt{\frac{2 \bar{\alpha}_t}{1 - \bar{\alpha}_t}} \cdot \E\left[ \norm{ \epsilon_\theta\left( x_t, t \right) - \epsilon }^2 \right] + \beta_t \norm{ \nabla_x \log p_t(x_t) }^2 + \frac{\beta_t}{2} \left\| \nabla_x \log p_t(x_t) - s_\theta(x_t, t) \right\|^2
\end{equation}

当信噪比趋于零时，前向过程收敛到标准正态分布。

\begin{equation}
\lim_{t \to T} q\left( x_t \mid x_0 \right) = \mathcal{N}\left( 0, I \right), \qquad \bar{\alpha}_t = \prod_{s=1}^{t} \left( 1 - \beta_s \right) \to 0
\end{equation}
\end{document}
